% Moved from the stationary section: numerical validation, kept out of the main line.
%======================================================================
\subsection{Implementation comparisons at finite wavelength}\label{app:validation}\label{sec:stationary-validation}
%======================================================================
The analytical statements above do not rest on numerics. To test the assembled representation
away from the gradient expansion, \path{solution/numerics/compare_cubic_quadrupole.py} compares
\eqref{eq:combined} with an independent forward calculation
(\path{solution/numerics/forward_cubic.py}) at $\bar t_\star=1.2$, $\bar t_o=3$, $\bar t_I=0.5$,
$|k|\simeq1$, so that $|k|\bar t_o\approx3$.
\begin{itemize}
\item \emph{Forward side.} The fields are solved through third order from (N) and the Cartesian (T)
  rows with the initial row. The constraint rows are not used. The Fermi--Walker triad is
  transported with $\dd E/\dd t=-SE$, and the photon geodesic is integrated by multilinear Picard
  iteration. The emission and clock roots and the exact normalisation follow. Time integrals are
  spectral and sky averages use Gauss quadrature.
\item \emph{Assembled side.} $J_{111}$ is the repository's second-ray-free direct measurement,
  with its nested integrals evaluated by quadrature. $S_{\rm meas}$ is
  \eqref{eq:measurement-source}, and $J_{\rm lin}$ is \eqref{eq:combined-linear}.
\end{itemize}
For three triples the two sides of \eqref{eq:combined} agree to $\le3\times10^{-9}$:
a noncoplanar triple with unequal frequencies, a triple with an exactly zero pair momentum, and
one with $|k_b+k_c|=0.05$. The pieces themselves are of order one to ten. In addition:
\begin{itemize}
\item $J_{111}$ agrees with the forward first-field-only map to $\le3\times10^{-8}$;
\item $J_{\rm lin}$ in constraint rows agrees with the Cartesian third-order field to $4\times10^{-14}$;
\item the forward map reproduces the exact channel \eqref{eq:dilation-channel} to $4\times10^{-8}$;
\item the no-probe-ray Hessian alone matches the forward mixed response to $10^{-9}$.
\end{itemize}
Two controls are resolved: replacing the Fermi--Walker triad by the symmetric one changes the
generic-triple value by $9\times10^{-5}$, and dropping the clock changes it by $0.24$.

\paragraph{What survives the observer dilation.} With the exact observer-centred identity of
\S\ref{sec:obs-dilation}, $\mathcal K_3=\mathcal K^c_3-\sum_{\rm pairs}\mathfrak D\mathcal K_2
-\sum_a\mathfrak D^2\mathcal K_1$, the same code isolates the genuinely cubic observer-centred kernel.
For the noncoplanar triple it is $0.9\%$ of $\mathcal K_3$: $+0.063$ of $-6.759$. For the
configuration $(k,q,-q)$ that enters $C^{(13)}$ it is $33\%$: $-2.29$ of $-6.89$. In these two
samples the dilation of the linear and quadratic maps carries most of the cubic response,
though the centred remainder is substantial in the second. Prop.~\ref{prop:obs-dilation} converts that dilation
part into $\frac12\sigma^2D^2C^{(11)}$ and a total derivative. The shell values are in
the exploratory spectra below.

\paragraph{Gradient coefficients.}
The forward calculation reproduces the degree-two and degree-four coefficients of
\eqref{eq:stationary-four-gradient} to $1.3\times10^{-9}$ and $1.6\times10^{-7}$,
at $\iota=0.5$ and $1$. Their difference reproduces the observer-local term to $10^{-6}$
in the script \path{solution/numerics/check_kappa4_shell.py}.
This is supplementary to the exact curvature and stationary calculations. At $x_o=0.3$, the leading centred coefficient gives
$-0.194$ in units of $\mathcal K_1$, against the numerical value $-0.195$; higher gradients
remain in the latter.

\paragraph{Top-order adjoint and linear-asymptotic comparisons.}
In \path{solution/numerics/explore_adjoint.py}, the momentum-source form of
Prop.~\ref{prop:adjoint} agrees with the direct cubic Green response to a few parts in
$10^6$ at $x_o=1,3$. The equivalent contact form, which differentiates the Hamiltonian
source numerically, agrees to $10^{-3}$--$10^{-4}$; the full map comparison at $x_o=1.5$
is limited by its grid to $4\times10^{-4}$. At $x_o=1$, $\mu=0.9$, the contact form has
large cancelling homogeneous, lapse and Hamiltonian contributions, approximately
$14$, $-73$ and $58$ in units of $T_\ell$, with an unrounded sum of $-0.24$.
These are implementation diagnostics, not additional analytical cancellations.
For the linear subhorizon transfer at $\ell=2$, $R=3$, the expansion through $O(x_o^{-3})$
has relative differences $1.8\times10^{-2}$, $3.5\times10^{-4}$ and $7.6\times10^{-8}$
from the exact transfer at $x_o=40,120,360$, respectively.

\subsection{Higher-gradient projection fits}\label{app:projection-fits}
The exact centred-$22$ proportionality of \S\ref{sec:obs-dilation} is established through
spectral degree six. Numerical fits suggest that it does not extend to degree eight:
\path{solution/numerics/centred22_ratio.py} gives ratios $52.66$ at $(e,o)=(1.2,3)$ and
$52.29$ at $(0.8,3.5)$, stable to $\le0.015$ under a change of fit points. The lower-order
ratios $5/3$ and $28/3$ are reproduced to $3\times10^{-6}$.
With the fixed template $\TF(\hat k\hat q+\hat q\hat k)$,
\path{solution/numerics/half_dilation_projection.py} gives relative mismatches from the
$-\mathfrak D\delta\mathcal K_1/2$ relation of $1.3\times10^{-4}$ at $k_0=0.1$ and
$2.4\times10^{-3}$ at $0.2$, suggesting a higher-order departure.
Neither fit proves the next analytical coefficient or an all-orders failure. A degree-eight
spectrum calculation needs both the four-gradient square and the two-/six-gradient
interference; that calculation remains open.

